% Generated from docs/REPORT_DRAFT.md; edit the manuscript then rerun the converter.
\chapter{Ứng dụng}

\section{Hỏi–đáp}
\reportfigure{qa_sequence.pdf}{Luồng xử lý một yêu cầu hỏi--đáp}{fig:qa-sequence}

\begin{Verbatim}[fontsize=\small,breaklines=true,frame=single]
Question → Qwen QueryPlan hoặc parser fallback → entity linking
         → whitelist compiler/catalog → parameterized Cypher
         → Neo4j → answer + evidence + latency
\end{Verbatim}

QueryPlan giới hạn operation, target, entity type, filter field/operator, sort và
limit bằng Pydantic JSON Schema. Qwen không sinh Cypher và không trả lời bằng kiến
thức riêng. Entity linker canonicalize tên trước query và trả confidence trong
evidence. Compiler chỉ phát sinh graph pattern đã whitelist; user input luôn nằm
trong parameter. Khi không có LLM, parser deterministic hỗ trợ chín intent.

\section{Recommendation}
\reportfigure{recommendation_explanation.pdf}{Luồng xếp hạng và tạo giải thích gợi ý phim}{fig:recommendation-flow}

Neo4j tạo candidate bằng các phim chia sẻ ít nhất một director, actor, keyword,
genre hoặc studio. Mỗi feature được tính document frequency trong graph và đóng
góp IDF theo type weight. Query trả candidate, score và feature chung; Python chỉ
định dạng \texttt{Recommendation} và explanation. Tie-break theo title giúp kết quả xác
định. Endpoint nhận \texttt{movie\_id} thay vì title để tránh nhập nhằng.

\section{API và UI}
\reportfigure{web_ui.pdf}{Giao diện hỏi--đáp và gợi ý phim}{fig:web-ui}

\begin{longtable}{|>{\raggedright\arraybackslash}p{0.44\textwidth}|>{\raggedright\arraybackslash}p{0.44\textwidth}|}
\caption{Tổng hợp api và ui}\label{tab:8-1} \\
\hline
\textbf{Endpoint} & \textbf{Chức năng} \\ \hline
\endfirsthead\hline \textbf{Endpoint} & \textbf{Chức năng} \\ \hline\endhead
GET \texttt{/health} & kiểm tra Neo4j connectivity \\ \hline
GET \texttt{/stats} & số node theo label và tổng relationship \\ \hline
GET \texttt{/entities/search} & full-text rồi fallback parameterized search \\ \hline
POST \texttt{/ask} & QA answer, intent, evidence, query time \\ \hline
POST \texttt{/recommend} & top-K recommendation và explanation \\ \hline
\end{longtable}

UI có hai tab QA và recommendation, autocomplete phim, lịch sử hiển thị và state
loading/error/success. Lịch sử chỉ ở frontend; backend hiện stateless.

\section{An toàn và vận hành}

Secret nằm trong \texttt{.env}; raw data không commit. LLM endpoint bind localhost;
remote demo dùng SSH tunnel. Public API không có write endpoint. Hạn chế hiện tại
là chưa có authentication/rate limiting, nên cấu hình phù hợp demo local chứ
không phải public production deployment.

\section{Failure modes}

\begin{longtable}{|>{\raggedright\arraybackslash}p{0.44\textwidth}|>{\raggedright\arraybackslash}p{0.44\textwidth}|}
\caption{Tổng hợp failure modes}\label{tab:8-2} \\
\hline
\textbf{Tình huống} & \textbf{Hành vi} \\ \hline
\endfirsthead\hline \textbf{Tình huống} & \textbf{Hành vi} \\ \hline\endhead
LLM chưa cấu hình & dùng parser 9 intent \\ \hline
LLM timeout/JSON sai & fallback parser \\ \hline
confidence thấp/thiếu slot & trả clarification \\ \hline
entity không tồn tại & trả entity-not-found/không tìm thấy \\ \hline
movie ID recommend sai & HTTP 404 \\ \hline
Neo4j không sẵn sàng & health 503 \\ \hline
tên fuzzy nhập nhằng & không tự chọn tied candidate \\ \hline
\end{longtable}

Failure behavior là một phần correctness. Hệ QA an toàn nên thừa nhận không biết
hoặc yêu cầu làm rõ thay vì chạy query với entity đoán sai.

\section{Tính faithful của explanation}

QA evidence gồm canonical entity/link confidence và record/path trả từ Neo4j.
Recommendation explanation không được model sinh: nó tổng hợp từ shared feature
đã đóng góp score. Cách làm giảm độ tự nhiên so với NLG explanation \cite{colas2023knowledge} nhưng
có tính faithful cao trong phạm vi hệ thống: mỗi lý do tương ứng quan hệ tồn tại
trong graph và contribution xác định.
